\chapter{Mơ Mộng Trong Giờ}
\chapterintro{Thân ngồi trong lớp nhưng đầu thường đi xa hơn cả cửa sổ}{Mơ mộng là quyền của tuổi trẻ, chỉ đừng chọn đúng lúc thầy cô đang chốt ý.}

\begin{lucbatbox}{Lục bát: Bài giảng thành nhạc nền}
Thầy cô giảng rất chân thành\
Em nghe câu trước, câu sau hóa mờ\
Tâm tư đi dạo bất ngờ\
Lúc ra sân nắng, lúc chờ tương lai\
Đến khi bị hỏi một bài\
Em như vừa mới quay ngay về trần
\end{lucbatbox}

\begin{tutuyetbox}{Thất ngôn tứ tuyệt: Cửa sổ hấp dẫn}
Ngoài kia mây gió bình yên quá
Trong này bài giảng nối ý sâu
Mắt em lại chọn nhìn nơi khác
Ô cửa nhỏ hơn câu đang bàn
\end{tutuyetbox}

\begin{ghichu}
Mơ mộng trong giờ đáng khịa ở chỗ rất thơ, nhưng thường làm phần kiến thức đang giảng mất luôn cơ hội vào đầu.
\end{ghichu}

\ngancach

\begin{lucbatbox}{Lục bát: Mộng xa hơn đề}
Đề kia mới viết hai dòng\
Em đã nghĩ đến viễn cảnh năm sau\
Từ trường, từ lớp, từ nhau\
Rồi quay lại thấy bài đau nhói lòng\
Ước mơ thì rất mênh mông\
Chỉ mong hiện tại đừng không một hàng
\end{lucbatbox}

\begin{tutuyetbox}{Thất ngôn tứ tuyệt: Mơ đúng lúc}
Khi cần chú ý em lại mơ
Khi cần học thuộc lại chờ hứng
Tuổi trẻ quả nhiều cớ đỡ lưng
Đỡ cho thói mơ đúng giờ thôi
\end{tutuyetbox}